\documentclass[10pt]{article}
% Shared LaTeX style for MIT 18.06 exam revision notes.
% Compile topic files with Tectonic, XeLaTeX, or LuaLaTeX.

\usepackage[a4paper,margin=0.72in]{geometry}
\usepackage{fontspec}
\usepackage{amsmath,amssymb,mathtools}
\usepackage{booktabs,array,tabularx}
\usepackage{enumitem}
\usepackage{titlesec}
\usepackage{fancyhdr}
\usepackage{microtype}
\usepackage{xcolor}

\setmainfont{Times New Roman}
\setsansfont{Arial}

\setlength{\parindent}{0pt}
\setlength{\parskip}{3.6pt}
\setlength{\abovedisplayskip}{6pt}
\setlength{\belowdisplayskip}{6pt}
\setlength{\abovedisplayshortskip}{4pt}
\setlength{\belowdisplayshortskip}{4pt}
\renewcommand{\arraystretch}{1.08}
\setlength{\tabcolsep}{4.5pt}

\titleformat{\section}
  {\normalfont\large\bfseries}
  {\thesection}
  {0.55em}
  {}
\titlespacing*{\section}{0pt}{10pt}{3pt}

\titleformat{\subsection}
  {\normalfont\normalsize\bfseries}
  {\thesubsection}
  {0.5em}
  {}
\titlespacing*{\subsection}{0pt}{7pt}{2pt}

\setlist[itemize]{leftmargin=1.35em,itemsep=1pt,topsep=2pt,parsep=0pt}
\setlist[enumerate]{leftmargin=1.55em,itemsep=1pt,topsep=2pt,parsep=0pt}

\pagestyle{fancy}
\fancyhf{}
\fancyfoot[C]{\scriptsize\color{black!45}MIT 18.06 Linear Algebra Exam Notes --- \thepage}
\renewcommand{\headrulewidth}{0pt}
\renewcommand{\footrulewidth}{0pt}

\newcommand{\notesubtitle}[1]{%
  {\centering\itshape #1\par}
  \vspace{2pt}
}

\newcommand{\source}[1]{%
  {\centering\small #1\par}
  \vspace{6pt}
}

\newcommand{\labelnote}[2]{%
  \par\smallskip
  \noindent\textbf{#1.}\quad #2
  \par\smallskip
}

\newcommand{\tightlabel}[1]{%
  \par\smallskip
  \noindent\textbf{#1.}\quad
}

\newcolumntype{Y}{>{\raggedright\arraybackslash}X}
\newcolumntype{L}[1]{>{\raggedright\arraybackslash}p{#1}}

\newenvironment{notetable}
  {\small\begin{tabularx}{\linewidth}}
  {\end{tabularx}}


\begin{document}

\begin{center}
{\LARGE 3.2--3.3 Nullspace, Special Solutions, and Complete Solutions\par}
\end{center}
\notesubtitle{Concise exam revision notes for homogeneous and nonhomogeneous systems.}
\source{Based on handwritten MIT 18.06 revision notes.}

\section{Core Idea}

\labelnote{Core idea}{Row reduction separates a linear system into forced variables and free variables. Pivot variables are forced by the equations; free variables create the directions in the nullspace.}

The two central systems are:
\[
Ax=0
\]
and
\[
Ax=b.
\]
The first finds input directions that disappear under \(A\). The second asks whether the right side \(b\) can be built from the columns of \(A\).

\section{Big Picture}

Let \(A\) be an \(m\times n\) matrix.

\begin{center}
\small
\begin{tabularx}{\linewidth}{@{}L{0.21\linewidth}Y Y@{}}
\toprule
Question & Meaning & Output \\
\midrule
\(Ax=0\) & Which input vectors are sent to zero? & The nullspace \(N(A)\). \\
\(Ax=b\) & Can the columns of \(A\) produce this \(b\)? & Zero, one, or infinitely many solutions. \\
Row reduction & Reveals pivots and free variables. & The structure of all solutions. \\
\bottomrule
\end{tabularx}
\end{center}

The rank \(r\) is the number of pivots. The number of free variables is
\[
n-r.
\]

\labelnote{Exam trigger}{If a problem says ``find all solutions'', expect an answer of the form}
\[
x=x_p+x_n.
\]
Here \(x_p\) is one particular solution and \(x_n\) runs through the nullspace.

\section{Nullspace vs Column Space}

\begin{center}
\small
\begin{tabularx}{\linewidth}{@{}L{0.19\linewidth}Y L{0.18\linewidth}Y@{}}
\toprule
Space & Definition & Lives in & What it tells you \\
\midrule
Nullspace \(N(A)\) & All \(x\) such that \(Ax=0\). & \(\mathbb{R}^n\) & Directions that do not change \(Ax\). \\
Column space \(C(A)\) & All linear combinations of the columns of \(A\). & \(\mathbb{R}^m\) & Right sides \(b\) for which \(Ax=b\) is solvable. \\
\bottomrule
\end{tabularx}
\end{center}

Why the dimensions differ:
\begin{itemize}
  \item \(x\) has one entry per column of \(A\), so \(x\in\mathbb{R}^n\).
  \item \(Ax\) and \(b\) have one entry per row of \(A\), so they live in \(\mathbb{R}^m\).
  \item Therefore \(N(A)\subseteq\mathbb{R}^n\), while \(C(A)\subseteq\mathbb{R}^m\).
\end{itemize}

\labelnote{Common mistake}{Do not place the nullspace in \(\mathbb{R}^m\). The nullspace contains input vectors, so its vectors have length \(n\).}

\section{Free Variables and Special Solutions}

Free variables appear in columns without pivots. Each free variable gives one independent nullspace direction.

\subsection{Solving pattern}

\begin{enumerate}
  \item Row-reduce \(A\) to echelon or reduced echelon form.
  \item Mark pivot variables and free variables.
  \item Set one free variable to \(1\) and all other free variables to \(0\).
  \item Solve for the pivot variables.
  \item The resulting vector is one special solution.
  \item Repeat once for each free variable.
\end{enumerate}

If the free variables are \(x_3\) and \(x_4\), use
\[
(x_3,x_4)=(1,0)
\qquad\text{and}\qquad
(x_3,x_4)=(0,1).
\]
These choices produce two special solutions \(s_1\) and \(s_2\).

\labelnote{Core idea}{The special solutions form a basis for the nullspace:}
\[
N(A)=\operatorname{span}\{s_1,\ldots,s_{n-r}\}.
\]

\labelnote{Exam trigger}{If \(n>m\), then \(r\le m<n\). There must be at least one free variable, so \(Ax=0\) has a nonzero solution.}

\section{Complete Solution}

For \(Ax=b\), first check consistency. If the system is consistent, all solutions are
\[
x=x_p+x_n.
\]

\begin{center}
\small
\begin{tabularx}{\linewidth}{@{}L{0.23\linewidth}L{0.20\linewidth}Y Y@{}}
\toprule
Part & Equation & Role & How to find it \\
\midrule
Particular solution \(x_p\) & \(Ax_p=b\) & One fixed solution. & Often set free variables to \(0\). \\
Nullspace part \(x_n\) & \(Ax_n=0\) & Directions added to \(x_p\). & Combine special solutions. \\
\bottomrule
\end{tabularx}
\end{center}

If \(s_1,\ldots,s_k\) are special solutions, then
\[
x=x_p+c_1s_1+c_2s_2+\cdots+c_ks_k.
\]
The constants \(c_1,\ldots,c_k\) are arbitrary.

Why this works:
\[
A(x_p+x_n)=Ax_p+Ax_n=b+0=b.
\]
Adding a nullspace vector keeps the same right side.

\labelnote{Common mistake}{Do not multiply the particular solution by an arbitrary constant. The free constants belong only to nullspace directions.}

\section{Rank, Pivots, and Free Variables}

\begin{center}
\small
\begin{tabularx}{\linewidth}{@{}L{0.27\linewidth}Y Y@{}}
\toprule
Quantity & Meaning & Exam use \\
\midrule
\(r=\operatorname{rank}(A)\) & Number of pivots. & Counts independent columns. \\
Pivot variables & Variables solved from equations. & Forced by the system. \\
Free variables & Variables chosen freely. & Create special solutions. \\
\(n-r\) & Number of free variables. & Dimension of \(N(A)\). \\
\bottomrule
\end{tabularx}
\end{center}

The rank-nullity count used here is
\[
\dim N(A)=n-r.
\]
\labelnote{Exam memory}{Full column rank controls uniqueness. Full row rank controls existence. Square full rank gives both.}

For \(Ax=0\), row operations do not change the solution set, so it is enough to solve
\[
Rx=0.
\]
For \(Ax=b\), reduce the augmented matrix \([A\ b]\). A zero coefficient row with a nonzero augmented entry means no solution.

\section{Full Column Rank vs Full Row Rank}

\begin{center}
\small
\begin{tabularx}{\linewidth}{@{}L{0.24\linewidth}L{0.18\linewidth}Y Y@{}}
\toprule
Case & Condition & Nullspace & Behavior of \(Ax=b\) \\
\midrule
Square full rank & \(r=m=n\) & Only \(0\). & Exactly one solution for every \(b\). \\
Full column rank & \(r=n\le m\) & Only \(0\). & Zero or one solution. \\
Full row rank & \(r=m\le n\) & Dimension \(n-m\). & At least one solution for every \(b\). \\
Not full rank & \(r<m\) and \(r<n\) & Dimension \(n-r\). & Zero or infinitely many solutions. \\
\bottomrule
\end{tabularx}
\end{center}

\subsection{Full column rank}

Every column has a pivot.
\begin{itemize}
  \item No free variables.
  \item \(N(A)=\{0\}\).
  \item If a solution exists, it is unique.
  \item Some right sides \(b\) may still be outside the column space.
\end{itemize}

\subsection{Full row rank}

Every row has a pivot.
\begin{itemize}
  \item The column space is all of \(\mathbb{R}^m\).
  \item \(Ax=b\) is solvable for every \(b\in\mathbb{R}^m\).
  \item If \(n>m\), free variables make the solution infinite.
\end{itemize}

\labelnote{Common mistake}{Full column rank gives uniqueness, not automatic existence. Full row rank gives existence for every \(b\), not uniqueness.}

\section{Rank-One Matrix Example}

A rank-one matrix has the form
\[
A=uv^{T}.
\]
Then
\[
Ax=u(v^{T}x).
\]
If \(u\ne0\), then \(Ax=0\) is equivalent to
\[
v^{T}x=0.
\]
So the nullspace consists of vectors perpendicular to \(v\).

Example:
\[
A =
\begin{bmatrix}
1 & 3 & 10\\
2 & 6 & 20\\
3 & 9 & 30
\end{bmatrix}
=
\begin{bmatrix}
1\\
2\\
3
\end{bmatrix}
\begin{bmatrix}
1 & 3 & 10
\end{bmatrix}.
\]

The nullspace condition is
\[
x_1+3x_2+10x_3=0.
\]

\begin{center}
\small
\begin{tabularx}{0.72\linewidth}{@{}YY@{}}
\toprule
Free choice & Resulting special solution \\
\midrule
\(x_2=1,\ x_3=0\) & \((-3,1,0)\) \\
\(x_2=0,\ x_3=1\) & \((-10,0,1)\) \\
\bottomrule
\end{tabularx}
\end{center}

Therefore
\[
N(A)=\operatorname{span}\{(-3,1,0),(-10,0,1)\}.
\]

\labelnote{Common mistake}{For \(A=uv^{T}\), the nullspace condition is \(v^{T}x=0\), not \(u^{T}x=0\). The vector \(v\) matches the input dimension.}

\section{Common Mistakes}

\begin{center}
\small
\begin{tabularx}{\linewidth}{@{}L{0.38\linewidth}Y@{}}
\toprule
Mistake & Correct exam habit \\
\midrule
Putting \(N(A)\) in \(\mathbb{R}^m\). & \(N(A)\subseteq\mathbb{R}^n\), because \(x\) has \(n\) entries. \\
Counting rows to find special solutions. & Count free variables: \(n-r\). \\
Giving only \(x_p\) for \(Ax=b\). & Add the whole nullspace: \(x=x_p+x_n\). \\
Scaling \(x_p\) by a free constant. & Only special solutions get arbitrary constants. \\
Treating full column rank as automatic solvability. & It gives uniqueness if a solution exists. \\
Treating full row rank as uniqueness. & It gives existence for every \(b\); free variables may remain. \\
Ignoring inconsistent rows. & Check \([A\ b]\) before writing a solution. \\
\bottomrule
\end{tabularx}
\end{center}

\section{Final Exam Checklist}

\tightlabel{Checklist}
\begin{enumerate}
  \item Write the size: \(A\) is \(m\times n\).
  \item Row-reduce and find the rank \(r\).
  \item Mark pivot variables and free variables.
  \item For \(Ax=0\), create one special solution per free variable.
  \item Check that the number of special solutions is \(n-r\).
  \item For \(Ax=b\), reduce \([A\ b]\) and check consistency.
  \item Find one particular solution \(x_p\).
  \item Add the nullspace:
  \[
  x=x_p+c_1s_1+\cdots+c_ks_k.
  \]
  \item Check dimensions: \(x\in\mathbb{R}^n\), \(b\in\mathbb{R}^m\).
  \item State the outcome directly: no solution, one solution, or infinitely many.
\end{enumerate}

\end{document}
